\chapter{Implémentation logicielle}

\section{Organisation du dossier ARMAGE}

Le dossier \texttt{ARMAGE} contient l'ensemble du code source de
l'application de pilotage du banc de test, ainsi que les ressources et
outils nécessaires à son fonctionnement et à sa distribution. Chaque
sous-dossier regroupe les fichiers liés à une couche fonctionnelle précise
(matériel, traitement d'image, interface graphique, utilitaires), ce qui
permet de localiser rapidement le code concerné par une modification~: tout
ce qui touche la caméra se trouve dans \texttt{camera/}, tout ce qui touche
le moteur dans \texttt{motor/}.

\begin{figure}[H]
    \DTsetlength{0.2em}{1em}{0.2em}{0.4pt}{1.6pt}
    \dirtree{%
        .1 ARMAGE/.
        .2 main.py.
        .2 requirements.txt.
        .2 settings.json.
        .2 ARMAGE.spec.
        .2 build\_windows.bat.
        .2 build\_linux.sh.
        .2 GxIAPI.dll.
        .2 Mode\_emploi\_ARMAGE.pdf.
        .2 README.md.
        .2 app.log.
        .2 camera/.
        .3 camera\_controller.py.
        .2 motor/.
        .3 motor\_controller.py.
        .2 processing/.
        .3 image\_processor.py.
        .2 gui/.
        .3 main\_window.py.
        .3 main\_page.py.
        .3 module\_widget.py.
        .3 dialogs.py.
        .3 save\_dialog.py.
        .3 settings\_manager.py.
        .3 enums.py.
        .2 utils/.
        .3 config.py.
        .3 report\_generator.py.
        .3 user\_manual.py.
        .3 manual\_assets/.
        .2 gxipy/.
        .2 assets/.
        .2 captures/\DTcomment{généré}.
        .2 build/\DTcomment{généré}.
        .2 dist/\DTcomment{généré}.
        .2 \_\_pycache\_\_/\DTcomment{généré, plusieurs occurrences}.
    }
    \caption{Arborescence du dossier ARMAGE}
    \label{fig:armage-arborescence}
\end{figure}

\subsection{Fichiers à la racine}

Le point d'entrée est \texttt{main.py}~: il configure le logging, désactive
la mise en veille de l'ordinateur pendant la durée d'un test, force le thème
clair de l'interface Qt et lance la fenêtre principale. Les dépendances
Python sont listées dans \texttt{requirements.txt}, installables via
\texttt{pip install -r requirements.txt}. \texttt{settings.json} est un
fichier de configuration persistante (port moteur, index caméra, mouvements
personnalisés, nombre de modules), lu et écrit par
\texttt{gui/settings\_manager.py}.

\texttt{ARMAGE.spec}, avec les scripts \texttt{build\_windows.bat} et
\texttt{build\_linux.sh}, servent à empaqueter l'application en exécutable
via PyInstaller~; \texttt{GxIAPI.dll} (le SDK caméra Daheng) est embarquée
dans cet exécutable pour Windows. Enfin, \texttt{Mode\_emploi\_ARMAGE.pdf}
est généré par \texttt{utils/user\_manual.py}, \texttt{README.md} contient
la documentation technique du projet, et \texttt{app.log} est le journal
d'exécution.

\subsection{Dossiers de code source}

Chaque périphérique a son propre dossier, isolé de l'interface~:
\texttt{camera/} pilote la caméra Daheng via le SDK \texttt{gxipy}
(connexion, exposition, gain, capture), et \texttt{motor/} pilote le moteur
TMCM-1021 via le protocole TMCL et la bibliothèque \texttt{pytrinamic}
(connexion, remontage, dévidage, impulsions, limites de couple).
\texttt{processing/} compare les images successives (métrique MSE) pour
détecter l'arrêt réel du mouvement et en déduire la réserve de marche.

\texttt{gui/} contient l'interface graphique PySide6/Qt qui orchestre ces
trois couches~: la fenêtre principale, la page listant les modules actifs,
le widget de module individuel, les boîtes de dialogue, et la gestion des
paramètres persistants. \texttt{utils/} regroupe les fonctions transverses~:
\texttt{config.py} centralise les constantes du projet, tandis que
\texttt{report\_generator.py} et \texttt{user\_manual.py} génèrent
respectivement le rapport PDF de fin de test et le mode d'emploi, via
\texttt{fpdf2}.

Enfin, \texttt{gxipy/} est le SDK caméra fourni par Daheng (code tiers, non
modifié), et \texttt{assets/} regroupe les icônes de l'application.

\subsection{Dossiers générés automatiquement}

Les dossiers \texttt{captures/}, \texttt{build/}, \texttt{dist/} et
\texttt{\_\_pycache\_\_/} ne contiennent que du contenu généré
automatiquement et ne doivent pas être édités manuellement~:
\texttt{captures/} sert d'historique des essais (images et rapports PDF),
\texttt{build/} et \texttt{dist/} sont produits par PyInstaller lors de
l'empaquetage (ce dernier contenant l'exécutable final), et
\texttt{\_\_pycache\_\_/} contient le bytecode Python compilé.

\subsection{Logique de rangement}

Le dossier suit trois principes. D'abord, une séparation entre le matériel,
la logique et l'interface~: chaque périphérique physique (caméra, moteur) a
son propre module isolé, indépendant de l'interface graphique, ce qui
facilite le test en mode simulation (sans matériel branché) et la
maintenance. Ensuite, une centralisation de la configuration~: un seul
fichier (\texttt{utils/config.py}) regroupe les constantes ajustables, ce
qui évite de disperser des valeurs « magiques » dans le code. Enfin, une
séparation entre le code source et les artefacts générés~: les dossiers
\texttt{build/}, \texttt{dist/}, \texttt{captures/} et les
\texttt{\_\_pycache\_\_/} ne contiennent que du contenu généré
automatiquement (compilation, exécutable, résultats de tests) et pourraient
être supprimés puis régénérés sans perte de code.

% ─── 1. FIXATION ────────────────────────────────────────────────────────────

\section{Gestion du posage et sélection du mouvement}

\subsection{Saisie des informations de test}

Avant de lancer une mesure, l'opérateur renseigne dans l'interface quatre
informations obligatoires : la référence du test, le numéro de série de la
montre, le nom de l'opérateur et le type de mouvement horloger. Le bouton
de démarrage reste désactivé tant que les trois premiers champs textuels ne
sont pas remplis, empêchant tout lancement sans identification complète du
test. Ces informations sont ensuite intégrées au rapport PDF généré en fin
de mesure.

\subsection{Base de données des mouvements}

Les paramètres mécaniques de chaque mouvement sont centralisés dans une
base de données de configuration. À chaque référence de mouvement correspond
un ensemble de quatre paramètres :

\begin{itemize}
    \item $N_{\text{total}}$ : nombre de tours du barillet nécessaires pour
          atteindre la pleine charge du ressort moteur ;
    \item $r$ : rapport de transmission entre l'axe d'entraînement et le
          barillet ;
    \item $N_{24\text{h}}$ : nombre de tours de l'axe d'entraînement à
          dévider pour qu'il ne reste que 24~heures de réserve de marche ;
    \item $C_{\text{max}}$ : couple maximal du barillet (mN$\cdot$m),
          utilisé pour calculer le couple moteur à appliquer.
\end{itemize}

Le tableau~\ref{tab:mouvements} présente les valeurs configurées pour les
mouvements actuellement disponibles dans le système.

\begin{table}[h]
    \centering
    \caption{Paramètres des mouvements horlogers configurés dans le système}
    \label{tab:mouvements}
    \begin{tabular}{lcccc}
        \hline
        \textbf{Mouvement} &
        $N_{\text{total}}$ (tr) &
        $r$ &
        $N_{24\text{h}}$ (tr) &
        $C_{\text{max}}$ (mN$\cdot$m) \\
        \hline
        Mvt19    & 14{,}15 & 3{,}33 & 35{,}8  & 6{,}70 \\
        Mvt01    & 10{,}30 & 3{,}33 & 24{,}0  & 10{,}80 \\
        Mvt200   & 13{,}20 & 3{,}33 & 32{,}3  & 7{,}90 \\
        Mvt200.1 & 13{,}20 & 3{,}33 & 32{,}3  & 7{,}90 \\
        Mvt110   & 24{,}00 & 2{,}84 & 51{,}7  & 6{,}81 \\
        Mvt111   & 14{,}00 & 2{,}96 & 30{,}26 & 4{,}06 \\
        \hline
    \end{tabular}
\end{table}

\subsection{Chargement automatique des paramètres}

Lorsque l'opérateur sélectionne un mouvement dans la liste déroulante,
l'identifiant du mouvement est immédiatement enregistré dans l'état interne
du module. Les paramètres associés ne sont pas affichés à l'écran mais sont
chargés automatiquement au moment du lancement du test, à la demande des
étapes de remontage et de dévidage décrites à la section précédente.

Ce mécanisme de chargement à la demande garantit que les calculs de tours
et de couple sont toujours cohérents avec le mouvement sélectionné au
moment du démarrage, sans risque de décalage si l'opérateur modifie sa
sélection après une première tentative.

La liste des mouvements disponibles peut être étendue par l'opérateur
directement depuis l'interface, via la fenêtre de paramètres. Les mouvements
ainsi ajoutés sont persistés dans un fichier de configuration JSON et
rechargés automatiquement à chaque démarrage de l'application, au même
titre que les mouvements intégrés.


% ─── 2. ACCOUPLEMENT ET MOTORISATION ────────────────────────────────────────

\section{Séquence de remontage et de dévidage}
\label{sec:motorisation}

\subsection{Vue d'ensemble de la séquence}

Avant de démarrer la phase de mesure, le système exécute une séquence
automatisée en deux étapes successives : le remontage et le dévidage. Ces
étapes sont pilotées par un moteur pas à pas TMCM-1021 (Trinamic) connecté
au système via une liaison RS485. Elles sont précédées d'un accouplement
mécanique entre le moteur et la tige du mouvement, réalisé manuellement lors
du montage (section~\ref{sec:accouplement}).

\subsection{Accouplement}
\label{sec:accouplement}

L'accouplement a pour but d'engager mécaniquement le moteur avec le
mécanisme de remontoir de la montre avant d'appliquer un couple significatif,
en garantissant que le nombre de tours effectué par le moteur soit le même
que celui effectué par la tige du mouvement.

Une automatisation de cette étape a été tentée~: le moteur devait effectuer
une rotation horaire de 0{,}5 tour à vitesse réduite (10~tr/min) et à faible
couple ($\approx 6$~mN$\cdot$m), un couple censé être insuffisant pour
entraîner le barillet mais suffisant pour engager le mécanisme d'accouplement.
Ce couple avait été déterminé à partir de la plage de fonctionnement définie
au chapitre précédent (figure~\ref{fig:plages_couples}), calculée à partir du
couple \emph{maximal} du barillet plutôt que de sa valeur minimale. En
pratique, le couple minimal réellement nécessaire pour entraîner le mouvement
s'est révélé bien inférieur à cette estimation~: il n'existait donc pas de
palier de couple où le moteur tourne sans entraîner également le mouvement.
Cette automatisation a par conséquent dû être abandonnée~; le code
correspondant a été commenté plutôt que supprimé, et reste directement
réutilisable si cette approche devait être reprise, par exemple en
s'appuyant sur la bride glissante présente sur tous les mouvements du tableau
fabricant pour absorber le décalage résiduel.

L'accouplement entre la fourche en U (moteur) et la fourche en T (tige)
(section~\ref{sec:accouplement_fourche}) se fait donc désormais à la main
lors du montage, avant le lancement de la séquence automatisée.

\subsection{Étape 1 : remontage du barillet}

Le remontage consiste à remonter le ressort moteur (barillet) jusqu'à son niveau maximal.
Le moteur effectue une rotation horaire dont l'amplitude est définie à partir des caractéristiques mécaniques du
mouvement sélectionné :

\begin{equation}
    \theta_{\text{remontage}} = N_{\text{total}} \times r
    \label{eq:remontage}
\end{equation}

où $N_{\text{total}}$ est le nombre de tours du barillet nécessaires pour
atteindre la pleine charge et $r$ le rapport de transmission entre l'axe
d'entraînement et le barillet. Ces paramètres sont définis pour chaque
référence de mouvement dans la base de données de configuration du système.

Le couple appliqué lors du remontage est calculé dynamiquement à partir du
couple maximal du barillet et du rapport de transmission, puis plafonné à
une valeur de sécurité de 15~mN$\cdot$m afin de protéger le mouvement :

\begin{equation}
    C_{\text{moteur}} = \min\!\left(
        \frac{C_{\text{barillet}}}{r} + C_{\text{détente}},\;
        C_{\text{max}}
    \right)
    \label{eq:couple}
\end{equation}

où $C_{\text{détente}} = 5$~mN$\cdot$m est le couple résistif magnétique à
vaincre et $C_{\text{max}} = 15$~mN$\cdot$m la limite de sécurité.
La vitesse de rotation est fixée à 50~tr/min pour tous les mouvements.

Ce couple cible $C_{\text{moteur}}$ est ensuite traduit en un palier CS du
driver~: le logiciel sélectionne le palier le plus proche au-dessus de la
cible (fonction \texttt{movement\_torque\_to\_tmcl()}, fichier
\texttt{config.py}). Trois régimes sont ainsi couverts~:
\begin{itemize}
    \item Régime 1 (5 -- 9~mN$\cdot$m)~: le moteur tourne mais ne peut pas
          entraîner le mouvement (accouplement) ;
    \item Régime 2 (9 -- 15~mN$\cdot$m)~: couple suffisant pour entraîner le
          mouvement ;
    \item Mode vieillissement~: accès à l'ensemble des 32 paliers, y compris
          sous le seuil de détente.
\end{itemize}
Lorsqu'un couple réellement nul est requis (roue libre), le registre CS
n'est pas descendu en dessous de son plancher~: le logiciel bascule sur un
mécanisme distinct, la consigne \texttt{StandbyCurrent~=~0} combinée à un
délai de roue libre (\texttt{FreewheelingDelay}), qui coupe l'étage de
puissance plutôt que de réduire davantage le courant (fonction
\texttt{deactivate\_torque()}, fichier \texttt{motor\_functions.py}).

\subsection{Étape 2 : dévidage partiel du barillet}

À l'issue du remontage, une pause de 5~secondes est observée pour laisser
le mécanisme se stabiliser. Le moteur effectue ensuite une rotation
anti-horaire afin de consommer une partie de l'énergie stockée dans le
barillet, de manière à ce qu'il reste exactement 24~heures de réserve de
marche :

\begin{equation}
    \theta_{\text{dévidage}} = N_{24\text{h}}
    \label{eq:devidage}
\end{equation}

où $N_{24\text{h}}$ est le nombre de tours de l'axe d'entraînement à
dévider pour qu'il ne reste que 24~heures de réserve de marche.
Ce paramètre est défini pour chaque référence de mouvement dans la base
de données de configuration.

Le couple et la vitesse appliqués lors du dévidage sont identiques à ceux
du remontage.

\subsection{Profil de mouvement trapézoïdal}

Pour les deux phases de rotation (remontage et dévidage), le contrôleur
moteur applique un profil de vitesse trapézoïdal : le moteur accélère
linéairement jusqu'à sa vitesse nominale, maintient cette vitesse sur la
majeure partie du déplacement, puis décélère symétriquement avant de
s'arrêter. Ce profil permet de limiter les à-coups mécaniques transmis
au mouvement horloger tout en minimisant la durée totale de la séquence.

La durée théorique de chaque phase est estimée par le système selon :

\begin{equation}
    t_{\text{total}} =
    \begin{cases}
        2\,t_{\text{acc}} + t_{\text{cst}}
            & \text{si } d > 2\,d_{\text{acc}} \\[4pt]
        2\sqrt{d / a}
            & \text{sinon}
    \end{cases}
    \label{eq:trapeze}
\end{equation}

où $t_{\text{acc}} = v_{\text{max}} / a$ est la durée de la phase
d'accélération, $d_{\text{acc}} = \tfrac{1}{2}\,a\,t_{\text{acc}}^2$
la distance parcourue pendant cette phase, $t_{\text{cst}}$ la durée
à vitesse constante, $d$ le déplacement total et $a$ l'accélération.
Cette estimation est utilisée pour afficher la progression en temps réel
dans l'interface.

\subsection{Démarrage de la mesure}

À l'issue du dévidage, le moteur est mis hors tension. L'accouplement
reste en place et la tige du mouvement demeure solidaire de l'arbre
moteur. C'est le couple de détente du moteur (5~mN$\cdot$m) qui suffit à
maintenir la tige bloquée en rotation, empêchant le barillet de se dévider
spontanément pendant la durée du test. La phase de mesure démarre
immédiatement.

% ─── 3. ACQUISITION OPTIQUE ─────────────────────────────────────────────────

\section{Pilotage du système d'acquisition optique}

\subsection{Initialisation de la caméra}

La caméra utilisée est une Daheng VEN-505-36U3C, pilotée via le SDK
Galaxy (\texttt{gxipy}). Son initialisation est effectuée au démarrage
de l'application et configure les paramètres suivants :

\begin{itemize}
    \item Résolution : $2592 \times 1944$~pixels (résolution
          maximale du capteur), afin de disposer du niveau de détail
          le plus élevé possible pour la comparaison d'images ;
    \item Temps d'exposition : 15\,000~µs, réglé
          expérimentalement pour obtenir une image correctement exposée
          dans les conditions d'éclairage contrôlé du banc de mesure ;
    \item Gain : 10~dB, choisi en complément du temps
          d'exposition pour équilibrer luminosité et bruit numérique ;
    \item Balance des blancs : mode \textit{Once}, déclenchée
          une seule fois à l'initialisation. Ce mode permet au capteur
          d'effectuer une compensation automatique des couleurs sur la
          première image, puis de figer ces coefficients pour toutes
          les captures suivantes, garantissant ainsi une cohérence
          colorimétrique entre images successives.
\end{itemize}

En l'absence de la caméra ou du SDK, l'application bascule
automatiquement en mode simulation : des images de test synthétiques
sont générées, ce qui permet de développer et tester le logiciel sans
matériel connecté.

\subsection{Contrôle de l'éclairage par GPIO}

L'éclairage du banc est piloté directement via le GPIO de la caméra
(broche Line2), configurée en sortie numérique. Ce choix présente deux
avantages : il n'nécessite aucun câblage supplémentaire ni
microcontrôleur dédié, et il garantit une synchronisation exacte entre
l'allumage de l'éclairage et la séquence d'acquisition, puisque les
deux sont pilotés par la même interface logicielle.

L'éclairage est activé uniquement pendant la durée nécessaire à chaque
capture, puis immédiatement éteint. Cette stratégie réduit l'exposition
thermique du mouvement horloger à la chaleur dégagée par la source
lumineuse et évite tout éclairement parasite entre deux captures.

\subsection{Séquence d'acquisition et délais de stabilisation}

Chaque capture suit une séquence rigoureuse destinée à garantir la
qualité et la reproductibilité des images :

\begin{enumerate}
    \item Allumage de l'éclairage via le GPIO ;
    \item Attente de stabilisation de 500~ms, permettant à l'exposition
          automatique résiduelle de converger et à l'éclairage
          d'atteindre son régime permanent ;
    \item Vidage du buffer d'acquisition (\textit{flush\_queue}),
          afin d'éliminer les trames accumulées pendant le délai de
          stabilisation et de s'assurer que l'image capturée correspond
          bien à l'instant courant ;
    \item Acquisition de l'image brute depuis le flux de la caméra,
          avec jusqu'à 10 tentatives espacées de 50~ms en cas d'échec ;
    \item Conversion de l'image du format Raw SDK vers un tableau
          NumPy en espace colorimétrique BGR ;
    \item Extinction de l'éclairage.
\end{enumerate}

Le délai de stabilisation de 500~ms a été retenu comme compromis entre
la qualité d'image et la durée totale de la séquence de capture. Il
constitue le principal facteur limitant pour la fréquence minimale
d'acquisition.

\subsection{Horodatage et sauvegarde sur disque}

Immédiatement après l'acquisition, l'heure exacte de la prise de vue
est gravée directement dans l'image sous la forme d'un bandeau
semi-transparent en haut du cadre, affichant la date et l'heure au
format \texttt{JJ/MM/AAAA HH:MM:SS}. Cette information est ainsi
indissociable de l'image elle-même, quelle que soit la façon dont
elle est consultée ultérieurement.

L'image est ensuite sauvegardée sur disque au format JPEG selon la
convention de nommage suivante :

\begin{center}
    \texttt{module\{id\}\_\{AAAAMMJJ\}\_\{HHMMSS\}.jpg}
\end{center}

où \texttt{\{id\}} est l'identifiant numérique du module de test et
\texttt{\{AAAAMMJJ\}\_\{HHMMSS\}} l'horodatage de la capture. Ce
nommage permet un tri chronologique naturel des fichiers et l'extraction
automatique de l'heure de prise de vue à partir du nom de fichier seul,
sans métadonnées supplémentaires.

L'ensemble des images d'un test est regroupé dans un répertoire
\texttt{captures/} situé à la racine de l'application. Aucune
organisation par sous-dossier n'est appliquée : le préfixe
\texttt{module\{id\}} dans le nom de fichier suffit à distinguer les
captures issues de modules différents lorsque plusieurs mouvements
sont testés en parallèle.

\section{Détection de l'arrêt du mouvement}
\label{sec:detection}

\subsection{Principe général}

La détection de l'arrêt du mouvement repose sur la comparaison d'images successives.
Le système capture une photographie du cadran à intervalles réguliers, et la compare à la précédente.
Tant que le mouvement est en marche, les aiguilles se déplacent entre les deux captures et par conséquent, les images diffèrent.
Lorsque le mouvement s'arrête, les aiguilles restent immobiles et les images deviennent donc quasiment identiques.
Cette différence peut être calculée et quantifiée par une métrique mathématique appelée MSE (\textit{Mean Squared Error}, erreur quadratique moyenne).

\subsection{Métrique de comparaison : le MSE}

Le MSE est une mesure de la différence pixel par pixel entre deux images.
Pour deux images $I_1$ et $I_2$ de $N$ pixels, le MSE est défini par :

\begin{equation}
    \text{MSE} = \frac{1}{N} \sum_{i=1}^{N} \bigl(I_1(i) - I_2(i)\bigr)^2
    \label{eq:mse}
\end{equation}

Le programme effectue également une conversion préalable en niveaux de gris.
Cette dernière élimine les variations de teinte liées aux légères fluctuations d'éclairage, tout
en conservant suffisamment d'informations pour détecter le déplacement
des aiguilles.

Un MSE élevé indique que les deux images sont significativement différentes.
En revanche, un MSE proche de zéro indique que les images sont quasiment identiques.

\subsection{Séquence de mesure}

La phase de mesure démarre immédiatement après la phase de dévidage. Le
système effectue une première capture à $t = 0$, qui sert de référence
initiale. À chaque intervalle de capture, une nouvelle image est acquise et
comparée à la précédente selon la séquence suivante :

\begin{enumerate}
    \item Allumage de l'éclairage ;
    \item Attente de stabilisation de l'exposition ;
    \item Capture d'une image en résolution maximale ($2592 \times 1944$~px) ;
    \item Conversion en niveaux de gris et calcul du MSE avec l'image précédente selon l'équation~\eqref{eq:mse} ;
    \item Estampillage horodaté et sauvegarde de l'image ;
    \item Extinction de l'éclairage.
\end{enumerate}

Si le MSE calculé est inférieur au seuil de détection, l'arrêt est confirmé
et la réserve de marche est enregistrée. Sinon, l'image courante devient la
nouvelle référence pour la prochaine comparaison.

\subsection{Intervalle de capture : compromis entre précision et consommation}
\label{sec:intervalle_capture}

Le choix de l'intervalle entre deux captures résulte d'un compromis entre
deux exigences antagonistes. D'une part, un intervalle court améliore la
précision temporelle de la mesure : l'erreur maximale sur la réserve de
marche est directement égale à la durée d'un intervalle. D'autre part, un
intervalle court implique un allumage fréquent de l'éclairage et une
sollicitation régulière de la caméra et du processeur, ce qui augmente la
consommation du système et peut poser problème lors de tests de longue durée.

En pratique, des intervalles de l'ordre de 10 à 15~minutes constituent un
bon équilibre : pour une réserve de marche typique de 40 à 70~heures,
l'erreur de mesure reste inférieure à 0,4~\%, ce qui est largement suffisant
pour caractériser un mouvement horloger.

\subsection{Détermination du temps de réserve de marche}

L'arrêt du mouvement est détecté lorsque la comparaison entre l'image actuelle $I_n$ et l'image précédente $I_{n-1}$ donne un MSE inférieur au seuil fixé.
Cela signifie que les aiguilles occupent la même position à $t_n$ et à $t_{n-1}$, ce qui signifie que la montre s'est arrêtée dans l'intervalle $[t_{n-2},\, t_{n-1}]$.
Par conséquent, le seul instant pour lequel le système peut garantir avec certitude que le mouvement fonctionnait encore est $t_{n-2}$.

La réserve de marche mesurée est donc définie par :

\begin{equation}
    R = t_{n-2} - t_0
    \label{eq:reserve_marche}
\end{equation}

où $t_0$ est l'instant où la mesure a débuté.

\subsection{Choix et calibration du seuil de détection}

Le seuil de détection (paramètre \texttt{MOVEMENT\_DETECTION\_THRESHOLD})
est la valeur critique du MSE en dessous de laquelle deux images consécutives
sont considérées comme identiques. Sa calibration répond à deux exigences
contradictoires :

\begin{itemize}
    \item Seuil trop élevé : risque de faux positifs. Un déplacement
          d'aiguille faible, par exemple lorsque la trotteuse est proche de
          sa position de la capture précédente, pourrait produire un MSE
          inférieur au seuil et déclencher une détection d'arrêt erronée.

    \item Seuil trop faible : risque de faux négatifs. Le bruit
          numérique du capteur de la caméra et les micro-vibrations
          environnementales génèrent un MSE résiduel non nul, même lorsque
          la scène est parfaitement immobile.
\end{itemize}

La valeur retenue a été déterminée expérimentalement à partir de plusieurs
essais.

\subsection{Robustesse et limites}

La méthode est robuste aux variations d'exposition grâce à la
conversion en niveaux de gris et à l'éclairage contrôlé et reproductible.
Une limite a néanmoins été identifiée : des vibrations mécaniques extérieures
importantes, provoquées par exemple par un choc ou une machine voisine,
peuvent modifier légèrement la position relative du cadran dans le champ de
la caméra et élever artificiellement le MSE, masquant temporairement un arrêt
réel.
